
\documentclass{article}
\usepackage{colortbl}
\usepackage{makecell}
\usepackage{multirow}
\usepackage{supertabular}

\begin{document}

\newcounter{utterance}

\centering \large Interaction Transcript for game `codenames', experiment `concrete', episode 5 with reka{-}flash{-}3.1{-}t0.0.
\vspace{24pt}

{ \footnotesize  \setcounter{utterance}{1}
\setlength{\tabcolsep}{0pt}
\begin{supertabular}{c@{$\;$}|p{.15\linewidth}@{}p{.15\linewidth}p{.15\linewidth}p{.15\linewidth}p{.15\linewidth}p{.15\linewidth}}
    \# & $\;$A & \multicolumn{4}{c}{Game Master} & $\;\:$B\\
    \hline

    \theutterance \stepcounter{utterance}  
    & & & \multicolumn{2}{p{0.3\linewidth}}{
        \cellcolor[rgb]{0.9,0.9,0.9}{
            \makecell[{{p{\linewidth}}}]{
                \texttt{\tiny{[GM$|$GM]}}
                {'hidden': {'team': ['shawl', 'butterfly', 'porch', 'pepperoni', 'nightshirt', 'carton', 'caribou', 'walrus', 'kleenex'], 'innocent': ['ball', 'matchstick', 'photo', 'lamb', 'bed', 'nectarine', 'pitchfork'], 'opponent': ['flute', 'beefburger', 'chopstick', 'warthog', 'owl', 'lamp', 'popcorn', 'firetruck'], 'assassin': ['bluebird']}, 'revealed': {'team': {'team': [], 'innocent': [], 'opponent': [], 'assassin': []}, 'opponent': {'team': [], 'innocent': [], 'opponent': [], 'assassin': []}}}
            }
        }
    }
    & & \\ \\

    \theutterance \stepcounter{utterance}  
    & & \multicolumn{4}{p{0.6\linewidth}}{
        \cellcolor[rgb]{0.9,0.9,0.9}{
            \makecell[{{p{\linewidth}}}]{
                \texttt{\tiny{[P1$\langle$GM]}}
                \texttt{Provide one single{-}word clue that relates to some of your team's words from the list below. You can choose to how many and to which words your clue relates to.} \\
\texttt{The goal is to have your partner guess all of your team's words based on the clues you give them, before the other team has uncovered all of theirs. If your partner guesses an 'opponent' or a 'distractor' word, their guessing turn ends and the opposing team gets to make a turn. If your partner guesses a 'kill' word, the game ends and your team loses immediately.} \\
\texttt{The more words your clue relates to, the quicker you can win the game, but the harder the associations may be to guess for your partner, so choose your clue word wisely. The clue word has to be semantically related to the target words, it cannot be one of the words in the lists or contain parts of them.} \\
\\ 
\texttt{Always give your single{-}word clue and your comma{-}separated list of related target words in the following format and make your answers as short as possible, never include any other text than is required in this form:} \\
\\ 
\texttt{CLUE: <WORD>} \\
\texttt{TARGETS: <TARGETS>} \\
\\ 
\texttt{Your team words are: shawl, butterfly, porch, pepperoni, nightshirt, carton, caribou, walrus, kleenex.} \\
\texttt{Your opponent's team words are: flute, beefburger, chopstick, warthog, owl, lamp, popcorn, firetruck.} \\
\texttt{Distractor words are: ball, matchstick, photo, lamb, bed, nectarine, pitchfork.} \\
\texttt{The kill words are: bluebird.} \\
            }
        }
    }
    & \\ \\

    \theutterance \stepcounter{utterance}  
    & \multicolumn{4}{p{0.6\linewidth}}{
        \cellcolor[rgb]{0.9,0.9,0.9}{
            \makecell[{{p{\linewidth}}}]{
                \texttt{\tiny{[P1$\rangle$GM]}}
                \texttt{<reasoning>} \\
\texttt{Let's tackle this clue{-}giving problem. The user's team has words like shawl, butterfly, porch, pepperoni, nightshirt, carton, caribou, walrus, kleenex. I need to find a single{-}word clue that relates to as many of these as possible without being one of the words themselves or containing parts of them.} \\
\\ 
\texttt{First, let's look for common themes. Let's go through each word:} \\
\\ 
\texttt{Shawl: A type of clothing, maybe fabric{-}related.} \\
\texttt{Butterfly: An insect, maybe something with wings or colors.} \\
\texttt{Porch: Part of a house, outdoor area.} \\
\texttt{Pepperoni: Food, specifically a type of sausage.} \\
\texttt{Nightshirt: Clothing worn at night.} \\
\texttt{Carton: Container, often for liquids or food.} \\
\texttt{Caribou: A type of reindeer, animal.} \\
\texttt{Walrus: Another Arctic animal.} \\
\texttt{Kleenex: Brand of tissues, paper product.} \\
\\ 
\texttt{Looking for connections. Let's see:} \\
\\ 
\texttt{Shawl, nightshirt, Kleenex – all are fabric or clothing items? Wait, Kleenex is a tissue, which is paper. Maybe not. Alternatively, maybe something related to use. Shawl is worn, nightshirt is worn, Kleenex is used for blowing nose. Maybe "use" but that's too vague.} \\
\\ 
\texttt{Butterfly, caribou, walrus – all animals. But they are different types: butterfly is insect, caribou and walrus are mammals. Maybe "animals" but that's too broad. Maybe "Arctic" since caribou and walrus are Arctic animals, but butterfly isn't. Hmm.} \\
\\ 
\texttt{Porch, carton – porch is a structure, carton is a container. Not obvious. Pepperoni is food. Maybe "food" but pepperoni is the only one here. Wait, caribou and walrus are animals, but maybe "mammals"? Caribou is a deer, which is a mammal, walrus is a marine mammal. Butterfly is an insect. So not all.} \\
\\ 
\texttt{Wait, let's think of other angles. Maybe "animals" for caribou, walrus, butterfly? No, butterfly is an insect. Maybe "Arctic animals" for caribou and walrus, but butterfly isn't. Alternatively, "reindeer" – caribou is a type of reindeer. So caribou and maybe walrus? No, wal} \\
            }
        }
    }
    & & \\ \\

    \theutterance \stepcounter{utterance}  
    & & & \multicolumn{2}{p{0.3\linewidth}}{
        \cellcolor[rgb]{0.9,0.9,0.9}{
            \makecell[{{p{\linewidth}}}]{
                \texttt{\tiny{[GM$|$GM]}}
                {'player': 'Cluegiver', 'type': 'rambling error', 'utterance': '<reasoning>\nLet\'s tackle this clue-giving problem. The user\'s team has words like shawl, butterfly, porch, pepperoni, nightshirt, carton, caribou, walrus, kleenex. I need to find a single-word clue that relates to as many of these as possible without being one of the words themselves or containing parts of them.\n\nFirst, let\'s look for common themes. Let\'s go through each word:\n\nShawl: A type of clothing, maybe fabric-related.\nButterfly: An insect, maybe something with wings or colors.\nPorch: Part of a house, outdoor area.\nPepperoni: Food, specifically a type of sausage.\nNightshirt: Clothing worn at night.\nCarton: Container, often for liquids or food.\nCaribou: A type of reindeer, animal.\nWalrus: Another Arctic animal.\nKleenex: Brand of tissues, paper product.\n\nLooking for connections. Let\'s see:\n\nShawl, nightshirt, Kleenex – all are fabric or clothing items? Wait, Kleenex is a tissue, which is paper. Maybe not. Alternatively, maybe something related to use. Shawl is worn, nightshirt is worn, Kleenex is used for blowing nose. Maybe "use" but that\'s too vague.\n\nButterfly, caribou, walrus – all animals. But they are different types: butterfly is insect, caribou and walrus are mammals. Maybe "animals" but that\'s too broad. Maybe "Arctic" since caribou and walrus are Arctic animals, but butterfly isn\'t. Hmm.\n\nPorch, carton – porch is a structure, carton is a container. Not obvious. Pepperoni is food. Maybe "food" but pepperoni is the only one here. Wait, caribou and walrus are animals, but maybe "mammals"? Caribou is a deer, which is a mammal, walrus is a marine mammal. Butterfly is an insect. So not all.\n\nWait, let\'s think of other angles. Maybe "animals" for caribou, walrus, butterfly? No, butterfly is an insect. Maybe "Arctic animals" for caribou and walrus, but butterfly isn\'t. Alternatively, "reindeer" – caribou is a type of reindeer. So caribou and maybe walrus? No, wal', 'message': 'Your answer contained more than two lines, please only give one clue and your targets on two separate lines.'}
            }
        }
    }
    & & \\ \\

\end{supertabular}
}

\end{document}
